%======================================================================
\section{Hilbert spaces and Dirac notation}\label{sec:hilbert}
%======================================================================

\begin{objectives}
  You will compute inner products and norms, prove and use the Cauchy--Schwarz
  inequality, read and write Dirac notation, expand a vector in an orthonormal
  basis and interpret the squared coefficients, and build orthonormal bases with the
  Gram--Schmidt procedure.
\end{objectives}

A vector space has no notion of length or angle: we can add vectors, but we
cannot say whether two of them are close, or perpendicular. Quantum mechanics
needs both. The overlap between two states will be a complex number, its modulus
squared will be a probability, and ``perpendicular'' will mean ``perfectly
distinguishable''. The structure that provides them is the inner product.

\subsection{Inner products and norms}

\begin{definition}[Inner product]\label{def:inner}
  An \term{inner product} on a complex vector space $V$ is a map
  $\langle\cdot,\cdot\rangle:V\times V\to\C$ such that, for all $u,v,w\in V$ and $\lambda\in\C$:
  \begin{enumerate}
    \item[(i)] $\langle v,w\rangle=\overline{\langle w,v\rangle}$ (conjugate symmetry);
    \item[(ii)] $\langle v,v\rangle\ge0$, with $\langle v,v\rangle=0$ only for $v=0$ (positivity);
    \item[(iii)] $\langle u,v+\lambda w\rangle=\langle u,v\rangle+\lambda\langle u,w\rangle$ (linearity in the second argument).
  \end{enumerate}
  A vector space with an inner product is called an \term{inner product space}; in finite dimension, it is also a \term{Hilbert space}.
  The \term{norm} of $v$ is $\norm{v}=\sqrt{\langle v,v\rangle}$.
\end{definition}

\begin{remark}
  In infinite dimension, a Hilbert space is required to be \emph{complete}, which
  means that Cauchy sequences converge. Every finite-dimensional inner product space is
  complete automatically, and we work in finite dimension almost everywhere, so
  we will not mention completeness again.
\end{remark}

Axioms (i) and (iii) imply that the inner product is \emph{conjugate-linear} in the first argument:
$\langle\lambda v,w\rangle=\bar\lambda\langle v,w\rangle$. This is the convention of physicists; mathematicians often
put the conjugate on the second argument instead.

\begin{pitfall}
  Scalars come out of the first slot \emph{conjugated} and out of the second one unchanged.
  Forgetting this is the most common source of sign and phase errors.
\end{pitfall}

\paragraph{Examples.}
On $\C^n$ the standard inner product is $\langle v,w\rangle=\sum_i\overline{v_i}\,w_i$, so $\norm{v}^2=\sum_i\abs{v_i}^2$. On the
continuous functions $[0,1]\to\C$ we use $\langle f,g\rangle=\int_0^1\overline{f(x)}\,g(x)\dd x$. On $n\times n$ matrices,
$\langle A,B\rangle=\Tr(A^\dagger B)$, which is the standard inner product of the lists of entries.

\begin{example}\label{ex:isinner}
  Which of the following are inner products on $\C^2$? (a) $\bar v_1w_1+2\bar v_2w_2$; (b) $\bar v_1w_1-\bar v_2w_2$;
  (c) $v_1w_1+v_2w_2$; (d) $\overline{(v_1+v_2)}(w_1+w_2)$.
\end{example}
\begin{solution}
  (a) Yes. Conjugate symmetry and linearity are clear, and $f(v,v)=\abs{v_1}^2+2\abs{v_2}^2>0$ unless $v=0$.

  (b) No: positivity fails. For $v=(0,1)$ we get $f(v,v)=-1<0$.

  (c) No: there is no conjugate, so it is not conjugate symmetric. It also violates positivity:
  $f(v,v)=1+\ii^2=0$ for $v=(1,\ii)\neq0$.

  (d) No: $f(v,v)=\abs{v_1+v_2}^2$ vanishes for the non-zero vector $v=(1,-1)$.
\end{solution}

\begin{example}\label{ex:fn-inner}
  In the space of continuous functions $[0,1]\to\C$ compute $\langle 1,x\rangle$, $\norm1$ and $\norm x$.
\end{example}
\begin{solution}
  $\langle1,x\rangle=\int_0^1x\dd x=\tfrac12$, $\ \norm1^2=\int_0^11\dd x=1$, and $\norm x^2=\int_0^1x^2\dd x=\tfrac13$, so
  $\norm1=1$ and $\norm x=1/\sqrt3$. Note that $\tfrac12\le1\cdot\tfrac1{\sqrt3}\approx0.577$, which is no accident: it
  is the Cauchy--Schwarz inequality below.
\end{solution}

The next result is the most important inequality in the subject. It says that the
overlap of two vectors can never be larger than the product of their lengths.

\begin{theorem}[Cauchy--Schwarz]\label{thm:cs}
  For all $u,v$ in an inner product space, $\ \abs{\langle u,v\rangle}\le\norm u\,\norm v$, with
  equality if and only if $u$ and $v$ are linearly dependent.
\end{theorem}
\begin{proof}
  If $v=0$ both sides are $0$, so assume $v\ne0$. Split $u$ into a part along $v$ and a remainder
  (Figure~\ref{fig:cs}):
  \[
    u=\frac{\langle v,u\rangle}{\norm v^2}\,v+w,\qquad
    w:=u-\frac{\langle v,u\rangle}{\norm v^2}\,v .
  \]
  Then $w\perp v$, because $\langle v,w\rangle=\langle v,u\rangle-\frac{\langle v,u\rangle}{\norm v^2}\langle v,v\rangle=0$.
  Since the two pieces are orthogonal, the Pythagorean theorem gives
  \[
    \norm u^2=\frac{\abs{\langle v,u\rangle}^2}{\norm v^2}+\norm w^2\ \ge\ \frac{\abs{\langle v,u\rangle}^2}{\norm v^2}.
  \]
  Multiplying by $\norm v^2$ and taking square roots gives the inequality, using $\abs{\langle v,u\rangle}=\abs{\langle u,v\rangle}$.
  Equality holds exactly when $w=0$, that is, when $u$ is a multiple of $v$.
\end{proof}

\begin{figure}[htb]
  \centering
  \begin{tikzpicture}[>=stealth,scale=1.1]
    \draw[very thick,navy,->] (0,0) -- (4,0) node[right] {$v$};
    \draw[very thick,navy,->] (0,0) -- (2.6,1.8) node[above right] {$u$};
    \draw[very thick,figacc,->] (0,0) -- (2.6,0) node[midway,below] {$\dfrac{\langle v,u\rangle}{\norm v^2}\,v$};
    \draw[very thick,ink,->] (2.6,0) -- (2.6,1.8) node[midway,right] {$w$};
    \draw[ink] (2.4,0) -- (2.4,0.2) -- (2.6,0.2);
  \end{tikzpicture}
  \caption{Proof of Cauchy--Schwarz: $u$ splits into a component along $v$ and an orthogonal remainder $w$.
  The component cannot be longer than $u$ itself.}
  \label{fig:cs}
\end{figure}

\begin{corollary}[Triangle inequality]\label{cor:tri-norm}
  $\norm{u+v}\le\norm u+\norm v$. Hence $\norm{\cdot}$ is a genuine norm.
\end{corollary}
\begin{proof}
  Expanding, $\norm{u+v}^2=\norm u^2+\norm v^2+2\Re\langle u,v\rangle\le\norm u^2+\norm v^2+2\norm u\norm v=(\norm u+\norm v)^2$,
  using $\Re z\le\abs z$ and Cauchy--Schwarz. This is the same proof as in Theorem~\ref{thm:triangle}.
\end{proof}

For unit vectors, Cauchy--Schwarz says $0\le\abs{\langle u,v\rangle}\le1$: overlaps of
normalised states are always legitimate probability amplitudes. In a real space it
also lets us define the angle $\theta$ between two vectors by
$\cos\theta=\langle u,v\rangle/(\norm u\norm v)$, since the right-hand side lies in $[-1,1]$.

\subsection{Dirac notation}

Physicists write vectors in a notation due to Dirac that is compact, a little strange at first, and indispensable once you get used to it.

\begin{itemize}
  \item A vector is written as a \term{ket} $\ket\psi$. A ket is a column vector in a chosen basis.
  \item The corresponding \term{bra} $\bra\psi$ is the row vector $\ket\psi^\dagger$, the conjugate transpose.
  \item The \term{bracket} $\braket\varphi\psi$ is the inner product $\langle\varphi,\psi\rangle$: the bra $\bra\varphi$ eats the ket $\ket\psi$ and returns a number.
  \item The \term{outer product} $\ket\psi\!\bra\varphi$ is the operator $\ket\chi\mapsto\braket\varphi\chi\,\ket\psi$: a column times a row, a matrix.
\end{itemize}
In $\C^n$ this is nothing but matrix multiplication: $\braket\varphi\psi=\varphi^\dagger\psi$ is a row times a column, a $1\times1$ matrix,
while $\ket\psi\!\bra\varphi=\psi\varphi^\dagger$ is a column times a row, an $n\times n$ matrix. The rules are the ones inherited from
the inner product: $\overline{\braket\varphi\psi}=\braket\psi\varphi$, the bra of $\lambda\ket\psi$ is $\bar\lambda\bra\psi$, and $\norm{\ket\psi}^2=\braket\psi\psi$.
Throughout the book, \emph{a state is a normalised ket}, $\braket\psi\psi=1$.

A bra can be seen as a \emph{linear form}: the map $\ket w\mapsto\braket vw$ is linear, so $\bra v$ belongs to the space $V^*=L(V,\C)$ of
linear forms on $V$, the \term{dual space}. Every linear form on a finite-dimensional Hilbert space arises this way, and the correspondence
$\ket v\mapsto\bra v$ is conjugate-linear and one-to-one. That is why bras and kets are two faces of the same information.

\begin{example}\label{ex:kets}
  Let $\ket\psi=\frac1{\sqrt3}\ket0+\frac{1+\ii}{\sqrt3}\ket1$. Check that it is normalised, write $\bra\psi$, compute $\braket0\psi$,
  and find the matrix $\ket\psi\!\bra\psi$ with its trace.
\end{example}
\begin{solution}
  As a column, $\ket\psi=\frac1{\sqrt3}\begin{pmatrix}1\\1+\ii\end{pmatrix}$, and $\braket\psi\psi=\frac13(1+\abs{1+\ii}^2)=\frac13(1+2)=1$.
  The bra is the conjugate transpose, $\bra\psi=\frac1{\sqrt3}\begin{pmatrix}1&1-\ii\end{pmatrix}$, and
  $\braket0\psi=\frac1{\sqrt3}$. The outer product is the column times the row:
  \[
    \ket\psi\!\bra\psi=\frac13\begin{pmatrix}1\\1+\ii\end{pmatrix}\begin{pmatrix}1&1-\ii\end{pmatrix}
    =\frac13\begin{pmatrix}1&1-\ii\\1+\ii&2\end{pmatrix}.
  \]
  Its trace is $\tfrac13(1+2)=1$. The matrix is Hermitian, as every $\ket\psi\!\bra\psi$ is.
\end{solution}

\begin{example}\label{ex:diracsimplify}
  Simplify $\braket01$, $\braket+-$, $(\ket0\!\bra1)(\ket1\!\bra0)$, $(\ket0\!\bra1)^2$ and $(\ket+\!\bra+)\ket0$.
\end{example}
\begin{solution}
  $\braket01=0$ and $\braket+-=\tfrac12(\braket0 0-\braket01+\braket10-\braket11)=\tfrac12(1-0+0-1)=0$.
  Next, bras and kets in the middle combine: $(\ket0\!\bra1)(\ket1\!\bra0)=\ket0\braket11\bra0=\ket0\!\bra0$,
  while $(\ket0\!\bra1)^2=\ket0\braket10\bra1=0$. Finally $(\ket+\!\bra+)\ket0=\ket+\braket+0=\tfrac1{\sqrt2}\ket+$.
\end{solution}

\subsection{Orthonormal bases}

\begin{definition}[Orthogonality]\label{def:ortho}
  Two vectors $\ket u,\ket v$ are \term{orthogonal} if $\braket uv=0$. A family $\{\ket{e_i}\}$ is
  \term{orthonormal} if $\braket{e_i}{e_j}=\delta_{ij}$, which means $1$ if $i=j$ and $0$ otherwise.
  An orthonormal family that is also a basis is an \term{orthonormal basis} (ONB).
\end{definition}

Non-zero pairwise orthogonal vectors are automatically independent: if $\sum_i\lambda_i\ket{e_i}=0$, apply $\bra{e_j}$ to get
$\lambda_j\norm{e_j}^2=0$, so $\lambda_j=0$. Hence $n$ orthonormal vectors in an $n$-dimensional space are an ONB. The point of an ONB
is that coordinates become trivial to compute.

\begin{theorem}[Expansion and Parseval]\label{thm:parseval}
  Let $\{\ket{e_i}\}$ be an ONB of $\Hil$. For every $\ket\psi\in\Hil$,
  \[
    \ket\psi=\sum_i\braket{e_i}\psi\,\ket{e_i},\qquad
    \norm{\psi}^2=\sum_i\abs{\braket{e_i}\psi}^2 .
  \]
\end{theorem}
\begin{proof}
  Write $\ket\psi=\sum_j\lambda_j\ket{e_j}$. Applying $\bra{e_i}$ gives $\braket{e_i}\psi=\sum_j\lambda_j\delta_{ij}=\lambda_i$, which is the first claim.
  For the second, $\norm\psi^2=\braket\psi\psi=\sum_{i,j}\overline{\lambda_i}\lambda_j\braket{e_i}{e_j}=\sum_i\abs{\lambda_i}^2$.
\end{proof}

No linear system has to be solved: the coordinate along $\ket{e_i}$ is the overlap $\braket{e_i}\psi$. For a normalised state, Parseval
says that the numbers $\abs{\braket{e_i}\psi}^2$ add up to $1$; when the postulates of Chapter~\ref{ch:postulates} come into play, they will be
the probabilities of the outcomes of a measurement in that basis. Reading the first formula as $\ket\psi=\bigl(\sum_i\ket{e_i}\!\bra{e_i}\bigr)\ket\psi$ gives the
\term{completeness relation}
\[
  \sum_i\ket{e_i}\!\bra{e_i}=\Id,
\]
valid for every ONB, which lets us ``insert an identity'' anywhere in a computation. For instance, inserting it on both sides of an operator $A$ gives
\[
  A=\sum_{i,j}\ket{e_i}\bra{e_i}A\ket{e_j}\bra{e_j}=\sum_{i,j}A_{ij}\ket{e_i}\!\bra{e_j},\qquad A_{ij}=\bra{e_i}A\ket{e_j},
\]
so the matrix element $A_{ij}$ is the $i$th coordinate of $A\ket{e_j}$, in agreement with the rule ``columns are images'' of Section~\ref{sec:vector-spaces}.

\begin{example}\label{ex:parseval}
  Expand $\ket\psi=\frac1{\sqrt3}\bigl(\ket0+(1+\ii)\ket1\bigr)$ in the ONBs $\{\ket0,\ket1\}$, $\{\ket+,\ket-\}$ and $\{\ket{+\ii},\ket{-\ii}\}$, and check Parseval.
\end{example}
\begin{solution}
  \emph{Computational basis.} $\braket0\psi=\frac1{\sqrt3}$ and $\braket1\psi=\frac{1+\ii}{\sqrt3}$. Squared moduli: $\tfrac13$ and $\tfrac23$, summing to $1$.

  \emph{Basis $\ket\pm$.} The bras are $\bra\pm=\frac1{\sqrt2}\begin{pmatrix}1&\pm1\end{pmatrix}$, so
  \[
    \braket+\psi=\frac{1}{\sqrt6}\bigl(1+(1+\ii)\bigr)=\frac{2+\ii}{\sqrt6},\qquad
    \braket-\psi=\frac{1}{\sqrt6}\bigl(1-(1+\ii)\bigr)=\frac{-\ii}{\sqrt6}.
  \]
  Squared moduli: $\tfrac56$ and $\tfrac16$, summing to $1$.

  \emph{Basis $\ket{\pm\ii}$.} Here $\bra{\pm\ii}=\frac1{\sqrt2}\begin{pmatrix}1&\mp\ii\end{pmatrix}$ (the second entry is conjugated), so
  \[
    \braket{+\ii}\psi=\frac1{\sqrt6}\bigl(1-\ii(1+\ii)\bigr)=\frac{2-\ii}{\sqrt6},\qquad
    \braket{-\ii}\psi=\frac1{\sqrt6}\bigl(1+\ii(1+\ii)\bigr)=\frac{\ii}{\sqrt6},
  \]
  with squared moduli $\tfrac56$ and $\tfrac16$, again summing to $1$. The same state has three different sets of coordinates, one for each point of view.
\end{solution}

\begin{pitfall}
  The formulas of Theorem~\ref{thm:parseval} need an \emph{orthonormal} basis. In $\R^2$, for the independent but
  non-orthogonal basis $(1,0),(1,1)$, the vector $(3,5)$ is $-2(1,0)+5(1,1)$, whereas the overlaps with the basis vectors are $3$ and $8$.
\end{pitfall}

\subsection{Gram--Schmidt}

Do orthonormal bases always exist? Yes, and the proof is a recipe. Given any basis, we make it orthonormal one vector at a time,
subtracting at each step the components along the vectors already produced (Figure~\ref{fig:gs}).

\begin{theorem}[Gram--Schmidt]\label{thm:gs}
  Let $\ket{v_1},\dots,\ket{v_n}$ be a basis of $\Hil$. Define recursively
  \[
    \ket{u_k}=\ket{v_k}-\sum_{j<k}\braket{e_j}{v_k}\ket{e_j},\qquad \ket{e_k}=\frac{\ket{u_k}}{\norm{u_k}},\qquad k=1,\dots,n .
  \]
  Then $\{\ket{e_k}\}$ is an orthonormal basis, and $\Span\{e_1,\dots,e_m\}=\Span\{v_1,\dots,v_m\}$ for every $m$.
\end{theorem}
\begin{proof}
  We argue by induction on $k$. Suppose $e_1,\dots,e_{k-1}$ are orthonormal and span the same space as $v_1,\dots,v_{k-1}$.
  For $i<k$, $\braket{e_i}{u_k}=\braket{e_i}{v_k}-\sum_{j<k}\braket{e_j}{v_k}\delta_{ij}=0$, so $u_k$ is orthogonal to all previous $e_i$.
  Moreover $u_k\neq0$: otherwise $v_k$ would be a combination of $e_1,\dots,e_{k-1}$, hence of $v_1,\dots,v_{k-1}$, contradicting independence.
  So $e_k$ is well defined, of norm $1$, and $\Span\{e_1,\dots,e_k\}=\Span\{v_1,\dots,v_k\}$.
\end{proof}

In particular, every finite-dimensional Hilbert space has an orthonormal basis.

\begin{figure}[htb]
  \centering
  \begin{tikzpicture}[>=stealth,scale=1.1]
    \draw[very thick,navy,->] (0,0) -- (3,0) node[right] {$e_1=v_1/\norm{v_1}$};
    \draw[very thick,navy!60,->] (0,0) -- (1.6,2.0) node[above] {$v_2$};
    \draw[dashed,ink] (1.6,2.0) -- (1.6,0);
    \draw[very thick,figacc,->] (0,0) -- (1.6,0) node[midway,below] {$\langle e_1,v_2\rangle e_1$};
    \draw[very thick,figacc,->] (1.6,0) -- (1.6,2.0) node[midway,right] {$u_2$};
    \draw[ink] (1.4,0) -- (1.4,0.2) -- (1.6,0.2);
  \end{tikzpicture}
  \caption{Gram--Schmidt, second step: removing from $v_2$ its component along $e_1$ leaves $u_2\perp e_1$,
  which after normalisation is $e_2$.}
  \label{fig:gs}
\end{figure}

\begin{example}\label{ex:gs-c2}
  Apply Gram--Schmidt to $\ket{v_1}=\ket0+\ket1$, $\ket{v_2}=\ket0$ in $\C^2$, and then in the reverse order.
\end{example}
\begin{solution}
  Since $\norm{v_1}=\sqrt2$, $\ket{e_1}=\frac1{\sqrt2}(\ket0+\ket1)=\ket+$. Then $\braket{e_1}{v_2}=\braket+0=\tfrac1{\sqrt2}$, and
  \[
    \ket{u_2}=\ket0-\tfrac1{\sqrt2}\ket+=\ket0-\tfrac12(\ket0+\ket1)=\tfrac12(\ket0-\ket1),\qquad
    \ket{e_2}=\frac{\ket{u_2}}{\norm{u_2}}=\ket-,
  \]
  because $\norm{u_2}=\tfrac1{\sqrt2}$. The result is the basis $\{\ket+,\ket-\}$.

  In the reverse order, $\ket{e_1}=\ket0$ and $\ket{u_2}=(\ket0+\ket1)-\braket0{v_1}\ket0=\ket1$, so $\ket{e_2}=\ket1$:
  now we obtain $\{\ket0,\ket1\}$. \emph{The outcome depends on the order of the input.}
\end{solution}

\begin{example}\label{ex:gs-r3}
  Apply Gram--Schmidt to $v_1=(1,1,0)$, $v_2=(1,0,1)$, $v_3=(0,1,1)$ in $\R^3$.
\end{example}
\begin{solution}
  $e_1=\frac1{\sqrt2}(1,1,0)$. Since $\langle e_1,v_2\rangle=\tfrac1{\sqrt2}$, we get
  $u_2=v_2-\tfrac1{\sqrt2}e_1=(1,0,1)-(\tfrac12,\tfrac12,0)=(\tfrac12,-\tfrac12,1)$, with $\norm{u_2}^2=\tfrac14+\tfrac14+1=\tfrac32$, so
  $e_2=\frac1{\sqrt6}(1,-1,2)$.
  For the third vector, $\langle e_1,v_3\rangle=\tfrac1{\sqrt2}$ and $\langle e_2,v_3\rangle=\tfrac1{\sqrt6}$, hence
  \[
    u_3=(0,1,1)-(\tfrac12,\tfrac12,0)-\tfrac16(1,-1,2)=(-\tfrac23,\tfrac23,\tfrac23),\qquad
    e_3=\tfrac1{\sqrt3}(-1,1,1).
  \]
  \emph{Check:} $e_1\cdot e_3=0$, $e_2\cdot e_3=\tfrac1{\sqrt{18}}(-1-1+2)=0$.
\end{solution}

\begin{exerciseset}
  \exercise{1} From axioms (i)--(iii) prove that $\langle\lambda v,w\rangle=\bar\lambda\langle v,w\rangle$, that $\langle v+v',w\rangle=\langle v,w\rangle+\langle v',w\rangle$, that $\langle0,w\rangle=0$,
  and that if $\langle v,w\rangle=0$ for all $w$ then $v=0$.

  \exercise{1} Decide whether each is an inner product on $\C^2$: $\bar v_1w_1+\bar v_2w_2$; $\bar v_1w_1-\bar v_2w_2$; $\bar v_1w_1$; $\bar v_1w_1+\bar v_2w_2+\bar v_1w_2+\bar v_2w_1$.

  \exercise{2} Let $A\in M_n(\C)$ and $f(v,w)=v^\dagger Aw$. Show that axiom (i) holds iff $A$ is Hermitian, and that $f$ is an inner product iff
  $A$ is Hermitian and $v^\dagger Av>0$ for all $v\neq0$. For which $c\in\C$ is $A=\begin{pmatrix}1&c\\\bar c&1\end{pmatrix}$ an inner product matrix?

  \exercise{2} On $C([0,1],\C)$ let $\langle f,g\rangle=\int_0^1\overline f\,g\dd x$. Check that this is an inner product (where do you use continuity?),
  compute $\langle1,x\rangle$ and $\norm x$, and show that $\{e^{2\pi\ii nx}\}_{n\in\Z}$ is an orthonormal system.

  \exercise{2} On $M_n(\C)$ define $\langle A,B\rangle=\Tr(A^\dagger B)$. Show that it is an inner product with $\norm A^2=\sum_{ij}\abs{A_{ij}}^2$,
  and that $\{\tfrac1{\sqrt2}\sigma_\mu\}_{\mu=0}^3$ is an orthonormal basis of $M_2(\C)$.

  \exercise{2} Show that $\norm{x}_1=\sum_i\abs{x_i}$ and $\norm x_\infty=\max_i\abs{x_i}$ are norms on $\C^n$, and sketch the unit balls of $\norm\cdot_1$, $\norm\cdot_2$,
  $\norm\cdot_\infty$ in $\R^2$.

  \exercise{3} Prove the polarisation identity, $\langle u,v\rangle=\tfrac14\bigl(\norm{u+v}^2-\norm{u-v}^2-\ii\norm{u+\ii v}^2+\ii\norm{u-\ii v}^2\bigr)$, with our convention. Conclude that the norm determines the inner
  product. Prove also the parallelogram law $\norm{u+v}^2+\norm{u-v}^2=2\norm u^2+2\norm v^2$ and show that $\norm\cdot_1$ and $\norm\cdot_\infty$ on $\R^2$ violate it.

  \exercise{1} Let $\ket\psi=\frac1{\sqrt3}\ket0+\frac{1+\ii}{\sqrt3}\ket1$ and $\ket\varphi=\frac1{\sqrt2}(\ket0-\ii\ket1)$. Write them as columns and the bras as rows, compute $\braket\psi\psi$, $\braket\varphi\varphi$, $\braket\varphi\psi$,
  and check $\braket\varphi\psi=\overline{\braket\psi\varphi}$. Compute $\ket\psi\!\bra\varphi$, $\ket\psi\!\bra\psi$ and $\ket\varphi\!\bra\varphi$: which are Hermitian, and what are their ranks and traces?

  \exercise{1} Simplify: $\bra0\sx\ket1$, $\bra1\sy\ket0$, $\bra+\sz\ket-$, $(\ket0\!\bra1)(\ket1\!\bra0)$, $(\ket+\!\bra+)\ket0$, $\bra{+\ii}\sy\ket{+\ii}$.

  \exercise{2} Prove: $(\ket\psi\!\bra\varphi)^\dagger=\ket\varphi\!\bra\psi$; $(\ket\psi\!\bra\varphi)(\ket\chi\!\bra\xi)=\braket\varphi\chi\,\ket\psi\!\bra\xi$; and $(\lambda\ket\psi)^\dagger=\bar\lambda\bra\psi$.

  \exercise{1} Write $\sx,\sy,\sz,H$ as $\sum_{ij}A_{ij}\ket i\!\bra j$, and compute the entries $A_{ij}=\bra{i}A\ket{j}$ of $\sy$ in the basis $\{\ket+,\ket-\}$.

  \exercise{1} Check that $\sum_i\ket{e_i}\!\bra{e_i}=\Id$ for the three bases $\{\ket0,\ket1\}$, $\{\ket\pm\}$, $\{\ket{\pm\ii}\}$.

  \exercise{2} For $\ket\psi=\cos\frac\theta2\ket0+e^{\ii\phi}\sin\frac\theta2\ket1$, compute $\bra\psi\sx\ket\psi$, $\bra\psi\sy\ket\psi$ and $\bra\psi\sz\ket\psi$, and show that the squares add up to $1$.

  \exercise{2} Show that $\ket\psi$ and $e^{\ii\gamma}\ket\psi$ give the same value of $\bra\psi A\ket\psi$ for every operator $A$ and of $\abs{\braket\varphi\psi}^2$ for every $\varphi$. Then show that
  $\frac1{\sqrt2}(\ket0+\ket1)$ and $\frac1{\sqrt2}(\ket0-\ket1)$ are \emph{not} equivalent in this sense, by finding an operator that tells them apart.

  \exercise{1} Which pairs are orthogonal? $\ket0,\ket+$; $\ket+,\ket-$; $\ket{+\ii},\ket{-\ii}$; $(1,\ii)$ and $(\ii,1)$; $(1,\ii)$ and $(1,-\ii)$; $(1,1,\ii)$ and $(1,-1,0)$.

  \exercise{2} Find all unit vectors of $\C^2$ orthogonal to $\frac1{\sqrt3}\bigl(\ket0+\sqrt2\,e^{\ii\pi/3}\ket1\bigr)$.

  \exercise{2} Let $\{\ket{e_i}\}$ be an ONB. Prove $\ket\psi=\sum_i\braket{e_i}\psi\ket{e_i}$ and $\norm\psi^2=\sum_i\abs{\braket{e_i}\psi}^2$ (Parseval), and show with an example in $\R^2$ that both fail for a non-orthogonal basis.

  \exercise{2} Expand $\ket\psi=\frac1{\sqrt3}\ket0+\frac{1+\ii}{\sqrt3}\ket1$ in the three bases of Example~\ref{ex:parseval} (checking the computations there), and interpret the squared coefficients.

  \exercise{2} Prove Pythagoras: $\braket uv=0\Rightarrow\norm{u+v}^2=\norm u^2+\norm v^2$. Show that the converse holds over $\R$ but fails over $\C$ (try $v=\ii u$).

  \exercise{3} Prove Bessel's inequality: if $\{\ket{e_1},\dots,\ket{e_k}\}$ is orthonormal but not necessarily a basis, then $\sum_{i\le k}\abs{\braket{e_i}\psi}^2\le\norm\psi^2$, with equality iff $\ket\psi\in\Span\{e_i\}$.

  \exercise{3} Two ONBs $\{\ket{e_i}\},\{\ket{f_j}\}$ of $\C^n$ are \emph{mutually unbiased} if $\abs{\braket{e_i}{f_j}}^2=1/n$ for all $i,j$. Show that the eigenbases of $\sx,\sy,\sz$ are pairwise
  mutually unbiased. Interpret: a state from one basis gives completely random outcomes when measured in another.

  \exercise{2} Let $B=\{(1,0),(1,1)\}\subset\R^2$. Find the dual basis $\{v_1^*,v_2^*\}$ as row vectors (so that $v_i^*(v_j)=\delta_{ij}$) and use it to compute the coordinates of $(3,5)$ in $B$.

  \exercise{3} Show that the dual basis of $\{\ket0,\ket+\}$ in $\C^2$ does not consist of the bras $\bra0,\bra+$, and explain why for an orthonormal basis it does.

  \exercise{1} Redo Example~\ref{ex:gs-c2} in both orders. Then apply Gram--Schmidt to $(1,1,0),(1,0,1),(0,1,1)$ in $\R^3$, and to $(1,\ii,0),(0,1,1),(1,0,1)$ in $\C^3$ in both orders.

  \exercise{2} Apply Gram--Schmidt to $(1,\ii),(\ii,1)$ in $\C^2$. What goes wrong with $(1,\ii),(\ii,-1)$, and why?

  \exercise{2} Prove that $\ket{u_k}\perp\ket{e_j}$ for $j<k$, that $u_k\neq0$, and that $\Span\{e_1,\dots,e_m\}=\Span\{v_1,\dots,v_m\}$. Show that if the $v_i$ are already orthogonal, Gram--Schmidt only normalises them.

  \exercise{3} Apply Gram--Schmidt to $\{1,x,x^2,x^3\}$ with $\langle f,g\rangle=\int_{-1}^1fg\dd x$. Compare with the Legendre polynomials $1,\ x,\ \tfrac12(3x^2-1),\ \tfrac12(5x^3-3x)$.

  \exercise{2} Extend $\ket{e_1}=\tfrac1{\sqrt3}(1,1,1)$ to an orthonormal basis of $\C^3$ using Gram--Schmidt. Find a different completion using $\tfrac1{\sqrt3}(1,\omega,\omega^2)$ and $\tfrac1{\sqrt3}(1,\omega^2,\omega^4)$ with $\omega=e^{2\pi\ii/3}$.
\end{exerciseset}
